% !TeX program = pdflatex
% ============================================================================
%  Mini-projet « Gas liquefaction (LNG) »
%  Cycles L0 (Linde-Hampson) et L3 (Claude, turbo-détendeur)
% ----------------------------------------------------------------------------
%  Les blocs oranges « [ PLACEHOLDER : INSÉRER UNE CAPTURE D'ÉCRAN ICI ] »
%  sont à remplacer par les captures d'écran du notebook. Pour insérer une
%  capture, supprimer le bloc \placeholder{...}{...} correspondant et le
%  remplacer par :
%    \begin{figure}[H]
%      \centering
%      \includegraphics[width=0.85\linewidth]{figures/ma_capture.png}
%      \caption{Légende.}
%      \label{fig:ma_capture}
%    \end{figure}
%  (le package float est déjà chargé pour l'option [H]).
% ============================================================================
\documentclass[11pt,a4paper]{article}

\usepackage[utf8]{inputenc}
\usepackage[T1]{fontenc}
\usepackage{lmodern}
\usepackage[french]{babel}
\usepackage{amsmath,amssymb}
\usepackage[margin=2.4cm]{geometry}
\usepackage{graphicx}
\usepackage{booktabs}
\usepackage{float}
\usepackage{xcolor}
\usepackage{caption}
\usepackage[hidelinks]{hyperref}

% ---------------------------------------------------------------------------
%  Boîte « placeholder » pour les captures d'écran du notebook.
%  Usage : \placeholder{Titre court}{Description précise de la capture}
% ---------------------------------------------------------------------------
\definecolor{phbox}{RGB}{200,100,0}
\definecolor{phfill}{RGB}{255,242,224}
\newcommand{\placeholder}[2]{%
  \begin{center}
  \fcolorbox{phbox}{phfill}{%
    \begin{minipage}{0.92\linewidth}
      \centering
      \vspace{1.6em}%
      {\Large\bfseries\color{phbox}[ PLACEHOLDER : INSÉRER UNE CAPTURE D'ÉCRAN ICI ]}\\[0.7em]
      {\bfseries #1}\\[0.4em]
      {\small\itshape #2}%
      \vspace{1.6em}%
    \end{minipage}}%
  \end{center}}

\captionsetup{font=small,labelfont=bf}

% ---------------------------------------------------------------------------
%  Page de garde (tous les champs entre crochets sont à compléter)
% ---------------------------------------------------------------------------
\begin{document}

\begin{titlepage}
  \centering
  \vspace*{2.5cm}
  {\large [ Établissement / École ] \par}
  \vspace{0.4cm}
  {\large [ Cours / UE ] \par}
  \vspace{2.5cm}
  {\Huge\bfseries Liquéfaction du gaz naturel \par}
  \vspace{0.8cm}
  {\LARGE Du cycle de Linde-Hampson (L0) au cycle de Claude (L3) \par}
  \vspace{0.5cm}
  {\large Mini-projet « Gas liquefaction (LNG) » : conception et comparaison
   de cycles de liquéfaction \par}
  \vspace{2.5cm}
  {\large\bfseries [ Prénom NOM 1 ] \par}
  {\large\bfseries [ Prénom NOM 2 ] \par}
  \vfill
  {\large [ Date ] \par}
  \vspace{1cm}
\end{titlepage}

% ===========================================================================
\section*{Introduction}
% ===========================================================================

La liquéfaction du gaz naturel (LNG) consiste à refroidir le méthane en
dessous de son point d'ébullition atmosphérique ($\approx -162\,^\circ$C à
1~bar), ce qui réduit son volume d'un facteur $\approx 600$ et rend possible
le stockage et le transport longue distance à grande échelle. Cette opération
est énergétiquement coûteuse et repose sur deux effets : le refroidissement
par détente de Joule-Thomson (un gaz réel se refroidit lorsqu'il est laminé
d'une haute vers une basse pression) et la récupération, dans laquelle la
vapeur froide non liquéfiée prérefroidit le gaz haute pression entrant.

Ce rapport s'inscrit dans le mini-projet « Gas liquefaction (LNG) », dont le
périmètre couvert ici est double : le cycle de base commun \textbf{L0}
(cycle de Linde-Hampson, étudié par tous les groupes) et le cycle avancé
\textbf{L3} (cycle de Claude avec turbo-détendeur). Les deux cycles sont
résolus en \emph{mode conception}, c'est-à-dire à un unique point de
fonctionnement en régime permanent, avec des propriétés de gaz réel
(CoolProp) et la bibliothèque partagée \texttt{caldis}. L'objectif est une
comparaison rigoureuse base/avancé à paramètres identiques, à l'aide des
métriques communes : rendement de liquéfaction, travail spécifique, figure
de mérite et surfaces d'échange.

% ===========================================================================
\section{Le cycle de référence L0 (Linde-Hampson)}
% ===========================================================================

\subsection{Architecture du système et hypothèses de modélisation}

Le cycle L0, schématisé ci-dessous, est le cycle de liquéfaction
récupératif le plus simple. Son fonctionnement est le suivant :
\begin{enumerate}
  \item le gaz d'appoint (méthane, $25\,^\circ$C, 1~bar) est mélangé à la
        vapeur recyclée dans un mélangeur ;
  \item le mélange est comprimé de 1 à 100~bar en quatre étages de rapport
        égal $r \approx 3{,}16$ ($1 \to 3{,}16 \to 10 \to 31{,}6 \to
        100$~bar), avec refroidissement à $T_0 + 5$~K $= 30\,^\circ$C entre
        chaque étage (intercoolers) et après le dernier étage (aftercooler),
        le rendement isentropique de chaque étage étant $\eta_{is} = 0{,}75$ ;
  \item le gaz haute pression est prérefroidi dans le récupérateur à
        contre-courant par la vapeur froide de retour (côté chaud : 100~bar ;
        côté froid : 1~bar ; coefficients de film $1000$ et $500$~W\,m$^{-2}$\,K$^{-1}$,
        ratio d'aires 1:1) ;
  \item il est ensuite laminé de 100 à 1~bar à travers la vanne
        Joule-Thomson (détente isenthalpique), ce qui produit un mélange
        diphasique ;
  \item le ballon de flash à 1~bar sépare le liquide saturé, soutiré comme
        produit ($\dot m_{liq} = 1$~kg/s, $\approx -162\,^\circ$C), de la
        vapeur saturée, qui traverse le récupérateur côté froid puis est
        recyclée à l'aspiration du compresseur.
\end{enumerate}

\placeholder{Schéma du cycle L0 (Linde-Hampson)}{Copier ici la figure
« The baseline cycle -- L0 » de l'énoncé (section 3), ou le schéma équivalent
dessiné dans le notebook. La flèche du gaz d'appoint, la boucle de vapeur
recyclée et les quatre étages de compression doivent être visibles.}

Les hypothèses de modélisation, communes à tous les cycles du projet, sont
les suivantes :
\begin{itemize}
  \item résolution en régime permanent, en mode conception uniquement (aucun
        mode simulation) ;
  \item propriétés de gaz réel (CoolProp) obligatoires : l'effet
        Joule-Thomson disparaît pour un gaz parfait, il est donc impossible
        de liquéfier un gaz idéal ;
  \item le récupérateur est modélisé par un échangeur \emph{mono-zone}
        (\texttt{SizingHX}) à $c_p$ constant ; le pincement de 5~K est imposé
        à un seul terminal, choisi selon le critère des taux de capacité
        (section~\ref{sec:pinch_L0}), et le minimum interne est vérifié
        \emph{a posteriori} ;
  \item seuls les récupérateurs sont dimensionnés en surface d'échange ; les
        inter/after-coolers sont de simples puissances thermiques ;
  \item bilan massique stationnaire : le débit d'appoint compense
        exactement le produit liquide, la vapeur non liquéfiée est recyclée.
\end{itemize}

Enfin, la boucle de recyclage fait du débit circulant une inconnue
implicite : elle n'est connue qu'une fois le rendement de liquéfaction $y$
résolu. Il s'agit d'un point fixe, initialisé avec $y_0 \approx 0{,}1$
(soit un débit circulant $\approx 10\,\dot m_{liq}$), sur lequel le solveur
converge.

\subsection{Grandeurs suivies et métriques de comparaison}

Les métriques sont évaluées au point de conception et définies comme suit.
\\
\\
Le rendement de liquéfaction est la fraction du débit circulant liquéfiée à
chaque passage :
\begin{equation}
  y = \frac{\dot m_{liq}}{\dot m_{circ}} .
  \label{eq:yield}
\end{equation}
La puissance de compression est la somme sur les quatre étages, chacun
caractérisé par son rendement isentropique
\begin{equation}
  \dot W_{comp} = \sum_{k=1}^{4} \dot m_{circ}\,\bigl(h_{k,s} - h_{k,e}\bigr),
  \qquad
  \eta_{is} = \frac{h_{2s} - h_1}{h_2 - h_1} = 0{,}75 ,
  \label{eq:comp}
\end{equation}
où $h_{2s}$ est l'enthalpie de sortie isentropique de l'étage. 
\\
\\
Le travail spécifique de liquéfaction, grandeur de référence du projet, est ramené au
kilogramme de LNG produit :
\begin{equation}
  w = \frac{\dot W_{tot}}{\dot m_{liq}},
  \qquad
  \dot W_{tot} = \dot W_{comp} \quad \text{(cycle L0)} .
  \label{eq:w}
\end{equation}
Il est comparé au travail minimal réversible de liquéfaction, obtenu en
intégrant le premier et le second principes entre l'état d'alimentation
($25\,^\circ$C, 1~bar) et le produit (liquide saturé, 1~bar), avec rejet de
chaleur à $T_0$ :
\begin{equation}
  w_{min} = T_0\,\bigl(s_{feed} - s_{liq}\bigr)
          - \bigl(h_{feed} - h_{liq}\bigr)
          \approx 1083~\text{kJ/kg} \quad \text{(méthane)} ,
  \label{eq:wmin}
\end{equation}
ce qui définit la figure de mérite
\begin{equation}
  \mathrm{FOM} = \frac{w_{min}}{w} .
  \label{eq:fom}
\end{equation}
S'y ajoutent, pour le récupérateur uniquement : la surface d'échange
($A$, $UA$) et la puissance échangée $\dot Q_{rec}$, issues de la
résolution en conception. Enfin, un contrôle qualité systématique est
effectué : la fermeture du bilan d'énergie du premier principe sur l'ensemble
du cycle doit être nulle à la précision du solveur.

\subsection{Le pincement du récupérateur : pourquoi l'imposer côté chaud}
\label{sec:pinch_L0}

Le récupérateur étant modélisé en une seule zone à $c_p$ constant, le
pincement de 5~K ne peut être imposé qu'à \emph{un seul} des deux terminaux
de l'échangeur à contre-courant : le bout chaud (entrée du flux chaud,
sortie du flux froid) ou le bout froid (sortie du flux chaud, entrée du flux
froid). Ce choix n'est pas libre : il est dicté par les taux de capacité des
deux flux.

Notons $C_h = (\dot m\, c_p)_h$ et $C_c = (\dot m\, c_p)_c$ les taux de
capacité des flux chaud (haute pression) et froid (basse pression). Le bilan
d'énergie global de l'échangeur et les écarts aux terminaux s'écrivent
\begin{align}
  \dot Q &= C_h\,\bigl(T_{h,e} - T_{h,s}\bigr)
         = C_c\,\bigl(T_{c,s} - T_{c,e}\bigr), \label{eq:hx_balance}\\
  \Delta T_{ch} &= T_{h,e} - T_{c,s},
  \qquad
  \Delta T_{fr} = T_{h,s} - T_{c,e}. \label{eq:hx_ends}
\end{align}
En soustrayant ces deux écarts et en utilisant \eqref{eq:hx_balance}, il
vient
\begin{equation}
  \Delta T_{ch} - \Delta T_{fr}
  = \dot Q \left( \frac{1}{C_h} - \frac{1}{C_c} \right)
  = \dot Q\,\frac{C_c - C_h}{C_h\, C_c} .
  \label{eq:pinch_criterion}
\end{equation}
Le signe de \eqref{eq:pinch_criterion} fixe le terminal porteur du pincement
sans aucune ambiguïté :
\begin{equation}
  C_h > C_c \iff \Delta T_{ch} < \Delta T_{fr}
  \quad\Longrightarrow\quad \text{pincement au bout chaud.}
  \label{eq:pinch_rule}
\end{equation}
L'interprétation physique est la suivante : le flux de plus petit taux de
capacité subit la plus grande excursion de température et « referme » l'écart
avec l'autre flux ; le minimum d'approche se situe donc au bout où ce flux
sort. Imposer le pincement au mauvais terminal produit un croisement des
profils de température, c'est-à-dire un design physiquement infaisable.
\\
\\
Pour le cycle L0, le critère \eqref{eq:pinch_rule} est nettement vérifié côté
chaud : le flux chaud transporte \emph{la totalité} du débit circulant à
100~bar (gaz dense, $c_p$ élevé), tandis que le flux froid ne transporte que
la vapeur non liquéfiée, soit $(1-y)\,\dot m_{circ} \approx
0{,}9\,\dot m_{circ}$ à 1~bar, avec un $c_p$ plus faible. On a donc $C_h >
C_c$ : le pincement de 5~K est imposé au bout chaud par la contrainte
\begin{equation}
  T_{h,e} - T_{c,s} = 5~\text{K} .
  \label{eq:pinch_constraint}
\end{equation}
Cette hypothèse est ensuite \emph{contrôlée a posteriori} : les profils de
température des deux flux sont reconstruits point à point le long de
l'échangeur (200 points, températures recalculées par CoolProp à partir des
enthalpies et du bilan d'énergie local), et le minimum interne
$\Delta T_{int} = \min (T_h - T_c)$ doit rester $\geq 5$~K. Pour L0, le
minimum du profil reste confondu avec le bout chaud, l'écart au bout froid
étant de l'ordre de 105~K : l'hypothèse est validée.

\subsection{Résultats du point de base}
\label{sec:resultats_L0}

\placeholder{Sortie « RÉSULTATS L0 » du notebook (cellule de calcul des
métriques)}{Copier la sortie console complète : rendement $y$, débit
circulant, $\dot W_{comp}$, $w$, $w_{min}$, FOM, pincements bout chaud/bout
froid, pincement interne minimal, fermeture du bilan d'énergie, surface
$A$, $UA$ et $\dot Q$ du récupérateur.}

Le point de base L0 confirme les ordres de grandeur attendus par l'énoncé
(les valeurs exactes figurent sur la capture ci-dessus) :
\begin{itemize}
  \item rendement de liquéfaction $y \approx 9\,\%$, soit un débit circulant
        $\dot m_{circ} \approx 10$ à 11~kg/s pour $\dot m_{liq} = 1$~kg/s ;
  \item travail spécifique $w \approx 1{,}1 \times 10^{4}$~kJ/kg
        ($\approx 3$ à 3{,}4~kWh par kg de LNG) ;
  \item figure de mérite $\mathrm{FOM} \approx 9$ à 10~\% (avec
        $w_{min} = 1083$~kJ/kg) ;
  \item surface du récupérateur, $UA$ et puissance échangée : voir la capture
        ci-dessous.
\end{itemize}
La fermeture du bilan d'énergie est vérifiée, et le pincement interne reste
au bout chaud, comme prévu à la section~\ref{sec:pinch_L0}. Ce point
constitue la référence à laquelle le cycle L3 sera comparé. On note dès
maintenant une contre-vérification utile : le cycle L3 évalué à
$x_e = 0{,}01$ (seul 1~\% du débit est dérivé vers le détendeur) est quasi
identique à L0 et retrouve exactement ces valeurs
(section~\ref{sec:resultats_L3}), ce qui conforte la cohérence des deux
modèles.

% ===========================================================================
\section{Le cycle L3 : ajout d'une turbine de détente (cycle de Claude)}
% ===========================================================================

\subsection{Modification apportée au cycle de base}

Le cycle L3, schématisé ci-dessous, reprend l'intégralité du cycle L0 et
y ajoute trois composants : un séparateur (\emph{splitter}) placé à la sortie
de l'aftercooler, une turbine de détente et un mélangeur en
amont du ballon de flash. Son fonctionnement est le suivant :
\begin{enumerate}
  \item à la sortie de l'aftercooler (30~$^\circ$C, 100~bar), une fraction
        $x_e$ du débit circulant est dérivée vers le turbo-détendeur, la
        fraction complémentaire $1 - x_e$ poursuivant le chemin classique
        (récupérateur puis vanne Joule-Thomson) ;
  \item le turbo-détendeur détend ce flux de 100~bar à 1~bar avec un
        rendement isentropique $\eta_{is} = 0{,}80$ ; parti de 30~$^\circ$C,
        le gaz ressort en vapeur monophasée légèrement surchauffée
        ($\approx -153\,^\circ$C soit $\approx 120$~K, à quelques kelvins de
        la saturation à 1~bar, $T_{sat} = 111{,}7$~K) : la turbine reste
        « sèche », sans aucune condensation interne ;
  \item la sortie froide du détendeur est mélangée au flux laminé juste
        avant le ballon de flash, dont la vapeur saturée rejoint, comme en
        L0, le côté froid du récupérateur.
\end{enumerate}
La fraction détendue est imposée par la contrainte de débit du splitter,
\begin{equation}
  \dot m_{exp} = x_e \, \dot m_{circ},
  \label{eq:split}
\end{equation}
et la puissance récupérée au détendeur, déduite du bilan enthalpique (indices
$in$/$out$ pour l'entrée et la sortie réelle, $out,s$ pour la sortie
isentropique),
\begin{equation}
  \dot W_{exp} = \dot m_{exp}\,\bigl(h_{in} - h_{out}\bigr),
  \qquad
  \eta_{is}^{exp} = \frac{h_{in} - h_{out}}{h_{in} - h_{out,s}} = 0{,}80 ,
  \label{eq:expander}
\end{equation}
vient en déduction du travail total, conformément à la convention du projet :
\begin{equation}
  \dot W_{tot} = \dot W_{comp} - \dot W_{exp} .
  \label{eq:w_L3}
\end{equation}
Le paramètre $x_e$ est la variable d'étude (valeur par défaut de l'énoncé :
$x_e = 0{,}5$) ; $x_e = 0$ redonne exactement le cycle L0.

\placeholder{Schéma du cycle L3 (Claude avec turbo-détendeur)}{Copier ici la
figure « L3 -- Claude cycle (with expander) » de l'énoncé (section 4), ou le
schéma équivalent. Le splitter après l'aftercooler, la turbine, le mélangeur
avant le flash et la vanne JT sur la branche principale doivent être
visibles.}

\paragraph{Mise en \oe uvre numérique.}
Conformément à la note numérique de l'énoncé, le cycle L3 ne peut pas être
résolu directement à $x_e = 0{,}5$ : la résolution utilise une continuation
en $x_e$. Le premier point $x_e = 0{,}01$ est initialisé à partir de la
solution L0 (les 97 variables communes aux deux cycles sont recopiées, les
24 nouvelles variables du splitter, du détendeur et du mélangeur étant
initialisées à partir de l'état L0) ; chaque point suivant est ensuite
résolu en \emph{warm-start} à partir du point déjà calculé le plus proche,
dont la solution complète est archivée. La boucle de recyclage reste un
point fixe, résolu à chaque $x_e$.

\subsection{Intérêt attendu de la modification}

L'ajout du turbo-détendeur présente \emph{a priori} trois bénéfices
cumulatifs :
\begin{enumerate}
  \item \textbf{Récupération de travail} : la puissance
        $\dot W_{exp}$ de \eqref{eq:expander} vient directement en
        déduction de la puissance de compression
        \eqref{eq:w_L3}, ce qui réduit le travail net.
  \item \textbf{Un refroidissement bien plus efficace que la détente
        laminée} : la vanne Joule-Thomson est isenthalpique (le gaz n'y perd
        aucune enthalpie) alors que la turbine extrait du travail ; pour une
        même chute de pression de 100 à 1~bar, l'enthalpie cédée au
        détendeur se traduit par une température de sortie
        bien plus basse. Depuis 30~$^\circ$C, la turbine livre une vapeur à
        $\approx 120$~K, alors que la branche laminée doit d'abord être
        prérefroidie par le récupérateur pour produire du liquide.
  \item \textbf{Un effet en cascade sur le rendement de liquéfaction} :
        l'injection de ce flux très froid dans le mélangeur situé avant le
        flash abaisse l'enthalpie à l'entrée du ballon ; la fraction
        liquéfiée $y$ augmente, donc le débit circulant
        $\dot m_{circ} = \dot m_{liq}/y$ nécessaire pour produire 1~kg/s de
        LNG diminue, et avec lui la puissance de compression. Le bénéfice
        sur $y$ se propage ainsi à toute la chaîne de compression.
\end{enumerate}
L'effet attendu, annoncé par l'énoncé, est donc un meilleur refroidissement,
un rendement de liquéfaction plus élevé, avec récupération de travail au
détendeur.

\subsection{Le côté du pincement devient un paramètre : méthode de choix}
\label{sec:choix_pinch}

En L3, l'hypothèse du pincement côté chaud, valable pour L0
(section~\ref{sec:pinch_L0}), ne peut plus être postulée une fois pour
toutes : le flux chaud du récupérateur ne transporte plus que la fraction
$(1 - x_e)\,\dot m_{circ}$, de sorte que le taux de capacité $C_h$
décroît lorsque $x_e$ augmente, tandis que le flux froid transporte toujours
la vapeur de retour $(1 - y)\,\dot m_{circ}$. Le rapport $C_h/C_c$ évolue
avec $x_e$ et le critère \eqref{eq:pinch_rule} doit être réévalué à chaque
point du balayage. La situation est en outre compliquée par le fait que le
$c_p$ du méthane varie fortement le long de l'échangeur (région
pseudo-critique à 100~bar) : le rapport des taux de capacité doit donc être
évalué \emph{à chaque terminal} (bout chaud et bout froid), et non
globalement.

La méthode mise en \oe uvre est une résolution à deux passes :
\begin{enumerate}
  \item \textbf{Première passe} : le pincement de 5~K est imposé côté chaud,
        comme en L0, pour chaque valeur de $x_e$ (contrainte
        \eqref{eq:pinch_constraint}) ; le profil interne de l'échangeur est
        ensuite reconstruit \emph{a posteriori} et son minimum
        $\Delta T_{int}$ comparé aux 5~K imposés.
  \item \textbf{Deuxième passe} : pour chaque $x_e$, le système est résolu
        deux fois, une fois avec le pincement imposé au bout chaud, une fois
        au bout froid (chaque côté disposant de sa propre chaîne de
        warm-starts) ; pour chaque candidat, le profil exact est recalculé
        et la faisabilité est définie par $\Delta T_{int} \geq 5$~K. Le
        design retenu est le candidat faisable ; si les deux le sont, on
        retient celui dont le minimum interne est le plus proche de 5~K,
        c'est-à-dire celui qui porte effectivement son pincement.
\end{enumerate}
Cette procédure garantit que le pincement imposé correspond au terminal
réellement critique, sans croisement de température ni design optimiste.

Un signal précurseur de la bascule est fourni par la première passe : l'écart
au bout froid, très confortable en L0 ($\approx 105$~K), se referme
progressivement lorsque $x_e$ augmente (tableau~\ref{tab:sweep}), signe que
le flux chaud, dont le débit diminue, perd de son avance en taux de capacité
et que le minimum d'approche migre progressivement du bout chaud vers
l'intérieur puis vers le bout froid de l'échangeur.

\subsection{Résultats du balayage en \texorpdfstring{$x_e$}{xe}}
\label{sec:resultats_L3}

Le tableau~\ref{tab:sweep} rassemble les résultats de la première passe
(pincement imposé au bout chaud) pour $x_e$ croissant de 0{,}01 à 0{,}50, à
produit liquide constant $\dot m_{liq} = 1$~kg/s.

\begin{table}[H]
  \centering
  \caption{Résultats du cycle L3 en fonction de la fraction détendue $x_e$
  (première passe : pincement de 5~K imposé au bout chaud ;
  $\dot m_{liq} = 1$~kg/s ; $\eta_{is}^{exp} = 0{,}80$). $\Delta T_{fr}$ est
  l'écart au bout froid et $\Delta T_{int}$ le minimum interne du profil
  reconstruit \emph{a posteriori}.}
  \label{tab:sweep}
  \small
  \setlength{\tabcolsep}{4pt}
  \begin{tabular}{ccccccccc}
    \toprule
    $x_e$ & $y$ & $\dot m_{circ}$ & $\dot W_{comp}$ & $\dot W_{exp}$ &
    $w$ & FOM & $\Delta T_{fr}$ & $\Delta T_{int}$ \\
     & (\%) & (kg/s) & (MW) & (MW) & (kJ/kg) & (\%) & (K) & (K) \\
    \midrule
    0,01 & 9,68  & 10,33 & 11,16 & 0,03 & 11\,132 & 9,7  & 105,1 & 5,00 \\
    0,05 & 10,97 & 9,11  & 9,85  & 0,13 & 9\,711  & 11,2 & 103,6 & 5,00 \\
    0,10 & 12,59 & 7,94  & 8,58  & 0,23 & 8\,346  & 13,0 & 101,6 & 5,00 \\
    0,15 & 14,21 & 7,04  & 7,60  & 0,31 & 7\,292  & 14,9 & 99,2  & 5,00 \\
    0,20 & 15,83 & 6,32  & 6,83  & 0,37 & 6\,453  & 16,8 & 96,5  & 5,00 \\
    0,30 & 19,06 & 5,25  & 5,67  & 0,46 & 5\,203  & 20,8 & 89,2  & 5,00 \\
    0,40 & 22,30 & 4,48  & 4,84  & 0,53 & 4\,316  & 25,1 & 78,1  & 5,00 \\
    0,50 & 25,54 & 3,92  & 4,23  & 0,58 & 3\,653  & 29,6 & 59,9  & 3,60$^{\dagger}$ \\
    \bottomrule
  \end{tabular}\\[0.2em]
  {\footnotesize $^{\dagger}$ pincement interne violé ($\Delta T_{int} <
  5$~K) : design infaisable en l'état, voir le texte.}
\end{table}

\placeholder{Tableau de synthèse du notebook (DataFrame du balayage
\texttt{xe})}{Copier le tableau du notebook affichant, pour chaque $x_e$ :
$y$, $\dot m_{circ}$, $\dot W_{comp}$, $\dot W_{exp}$, $\dot W_{tot}$, $w$,
FOM, $\Delta T_{int}$ (et la surface $A$ si disponible).}

Trois enseignements se dégagent.

\paragraph{Le point $x_e = 0{,}01$ valide le modèle.}
Avec seulement 1~\% du débit détendu, le cycle L3 est quasi identique à L0 :
$y = 9{,}68\,\%$, $w = 11\,132$~kJ/kg, $\mathrm{FOM} = 9{,}7\,\%$, soit les
valeurs du cycle de base à l'initialisation près : c'est la contre-vérification
annoncée à la section~\ref{sec:resultats_L0}.

\paragraph{L'effet attendu est confirmé et il est spectaculaire.}
Lorsque $x_e$ croît de 0{,}01 à 0{,}40, le rendement de liquéfaction passe de
9{,}7~\% à 22{,}3~\% (facteur $\approx 2{,}3$), le débit circulant chute de
10{,}3 à 4{,}5~kg/s, et le travail spécifique tombe de 11\,132 à
4\,316~kJ/kg, soit une réduction de 61~\% par rapport au point de base L0
(1,20~kWh/kg contre 3,09). La figure de mérite est multipliée par 2{,}6,
de 9{,}7~\% à 25{,}1~\%. La décomposition du bilan
(tableau~\ref{tab:sweep}) montre que la récupération de travail du
détendeur, bien que réelle ($\dot W_{exp}$ croît de 0{,}03 à 0{,}53~MW),
reste un effet du second ordre : au point $x_e = 0{,}4$, elle ne représente
qu'environ 11~\% de la puissance de compression. L'essentiel du gain provient
donc de l'effet en cascade décrit en section~2.2 : l'apport de froid de la
branche détendue au mélangeur du flash augmente $y$, ce qui réduit le débit
circulant et donc la puissance à comprimer.

\paragraph{La faisabilité du pincement limite le gain.}
Jusqu'à $x_e = 0{,}40$, le minimum interne du profil reste confondu avec le
pincement imposé au bout chaud ($\Delta T_{int} = 5{,}00$~K) : les designs
sont physiquement faisables. Au point $x_e = 0{,}50$, en revanche, la
reconstruction \emph{a posteriori} révèle un minimum interne de
$\Delta T_{int} = 3{,}60$~K $< 5$~K, tandis que l'écart au bout froid s'est
refermé à 59{,}9~K : le minimum d'approche a migré \emph{à l'intérieur} de
l'échangeur. Le design « pincement chaud » devient alors optimiste (le
modèle mono-zone à $c_p$ constant sous-estime l'écart réel près de la
région pseudo-critique du méthane), et le point $x_e = 0{,}50$ du
tableau~\ref{tab:sweep} doit être considéré comme infaisable en l'état.
Conformément à la méthode de la section~\ref{sec:choix_pinch}, la poursuite
du balayage au-delà de $x_e \approx 0{,}4$ doit donc être menée en imposant
le pincement au bout froid, le rapport des taux de capacité au bout chaud
étant alors passé sous l'unité.

\placeholder{Figure 2$\times$2 du notebook : $y$, $w$, FOM et
$\Delta T_{int}$ en fonction de $x_e$}{Copier le graphique en quatre
panneaux du notebook (tracé du DataFrame du balayage). Les points où le
pincement interne est violé y figurent entourés en rouge : vérifier que
seul $x_e = 0{,}5$ l'est.}

\placeholder{Diagrammes $T$--$Q$ et profils $\Delta T$ du récupérateur
(fonction \texttt{plot\_hx})}{Copier le tracé à gauche : profils $T_h$ et
$T_c$ en fonction de la chaleur cumulée à $x_e = 0{,}5$ ; et à droite :
$\Delta T(Q)$ le long du récupérateur pour plusieurs $x_e$ (0,3 / 0,4 /
0,5 / 0,6), avec les lignes de référence à 5~K et 0~K. Le creux sous 5~K
à $x_e = 0{,}5$ doit être visible.}

\placeholder{Résultats de la passe à deux côtés (tableau et figure de
bascule des taux de capacité)}{Copier : (i) le tableau du notebook avec, pour
chaque $x_e$, le côté retenu (chaud/froid), la faisabilité, $y$, $w$, FOM,
$\Delta T_{int}$ et les rapports $\dot m c_p$ aux deux bouts ; (ii) la figure
du rapport $\dot m c_p_{hot}/\dot m c_p_{cold}$ au bout chaud et au bout
froid en fonction de $x_e$, montrant la bascule sous l'unité.}

\subsection{Discussion}

\paragraph{Comparaison base / avancé.}
Au dernier point pleinement faisable de la première passe ($x_e = 0{,}40$),
le cycle de Claude divise par 2{,}6 le travail spécifique de liquéfaction par
rapport au Linde-Hampson ($4\,316$ contre $\approx 11\,100$~kJ/kg) et
multiplie la figure de mérite par 2{,}6 (25{,}1~\% contre 9{,}7~\%), pour un
rendement de liquéfaction lui aussi multiplié par 2{,}3. La comparaison
confirme l'intérêt du détendeur, avec la nuance quantitative importante que
la récupération de travail n'est pas le mécanisme dominant : c'est l'effet
thermique de la détente avec travail, via l'augmentation de $y$, qui porte
le gain.

\paragraph{Limites du modèle.}
Le modèle mono-zone à $c_p$ constant du récupérateur est une approximation
explicitement assumée (hypothèse C de l'énoncé), dont l'optimisme se
manifeste précisément pour L3, où les deux courbes $T(h)$ se resserrent :
la vérification \emph{a posteriori} du pincement interne n'est donc pas un
contrôle de forme mais une étape constitutive de la conception. Les valeurs
du point $x_e = 0{,}50$ (FOM $\approx 29{,}6\,\%$) sont à considérer comme
une borne optimiste tant que le design correspondant n'a pas été résolu avec
le pincement côté froid et vérifié.

\paragraph{Perspectives.}
Trois prolongements naturels : (i) un balayage fin entre $x_e = 0{,}40$ et
0{,}50 pour localiser précisément la bascule du pincement ; (ii) la
poursuite du balayage $x_e > 0{,}5$ avec la méthode à deux côtés, afin de
déterminer le $x_e$ optimal du cycle ; (iii) l'optimisation conjointe de
$x_e$ et de la haute pression $P_{high}$, tous deux cités par l'énoncé comme
paramètres d'optimisation.

% ===========================================================================
\section*{Conclusion}
% ===========================================================================

Ce rapport a modélisé et comparé, en mode conception et avec des propriétés
de gaz réel, le cycle de base L0 (Linde-Hampson) et le cycle avancé L3
(Claude avec turbo-détendeur) pour la liquéfaction d'un débit de 1~kg/s de
méthane. Le cycle L0, dont le pincement de récupérateur est imposé au bout
chaud en vertu du critère des taux de capacité
\eqref{eq:pinch_rule}, atteint les valeurs de référence attendues :
$y \approx 9\,\%$, $w \approx 1{,}1 \times 10^{4}$~kJ/kg, FOM
$\approx 9$ à 10~\%. L'ajout d'un turbo-détendeur dérivant une fraction
$x_e$ du débit haute pression améliore très nettement ces grandeurs : à
$x_e = 0{,}40$, le travail spécifique est réduit de 61~\% et la figure de
mérite atteint 25{,}1~\%, l'essentiel du gain provenant de l'augmentation du
rendement de liquéfaction plutôt que de la récupération de travail elle-même.
Le contrôle \emph{a posteriori} des profils de température du récupérateur
s'est révélé décisif : au-delà de $x_e \approx 0{,}4$, le pincement migre de
l'extrémité chaude vers l'intérieur de l'échangeur, rendant les designs
« pincement chaud » infaisables et imposant une méthode systématique de choix
du côté du pincement, appliquée sur l'ensemble du balayage.

\end{document}
